\documentclass[preprint]{sigplanconf}


%include agda.fmt

%subst keyword a = "\keyword{" a "}"
%subst conid a = "\identifier{" a "}"
%subst varid a = "\identifier{" a "}"
%format (HIDE(x)) = "{}^{\color{gray}" x "}"
%format (HIDE'(x)) = "{}_{\color{gray}" x "}"
%format = = "\kern.15em{=}\kern.15em"
%format rewrite = "\keyword{rewrite}"
%format -> = →
%format Set = "\constructor{Set}"
%format Set1 = "\constructor{Set_1}"
%format Set₁ = Set1
%format : = "{:}"
%format λ = "\lambda"
%format - = "\text-"
%format . = "\kern.6\underscorespace\text.\kern.6\underscorespace"
%format SEP = "\kern1.55\underscorespace;"
%format ABBRV = "\shade{\rlap{\ldots}\phantom{\hskip 1.05em}}"
%format COLSEP = "\qquad\qquad"
%format tt = "\constructor{tt}"
%format [] = "[\kern1pt]"
%format _≡_ = _ "\kern\underscorespace{\equiv}" _
%format ≡ = "\equiv"
%format refl = "\constructor{refl}"
%format ↦ = "\mathrel{\mapsto}"
%link ↦ _↦_
%format _↦_ = _ "\kern.5\underscorespace{" ↦ "}\kern-\underscorespace" _
%format MAPSTO = ↦
%format /↦ = "\mathrel{\not\mapsto}"
%format _/↦ = _ "\kern.5\underscorespace{" /↦ "}"
%format × = "{\times}"
%format _×_ = _ "{" × "}\kern-\underscorespace" _
%format _,_ = _ "\kern\underscorespace" , _
%format ⊎ = "\kern\underscorespace{\uplus}"
%format _⊎_ = _ "{" ⊎ "}" _
%format inj₁ = "\constructor{inj_1}"
%format inj₂ = "\constructor{inj_2}"
%format Maybe = "\constructor{Maybe}"
%format nothing = "\constructor{nothing}"
%format just = "\constructor{just}"
%format Bool = "\constructor{Bool}"
%format String = "\constructor{String}"
%format true = "\constructor{true}"
%format false = "\constructor{false}"
%format if = "\mathbf{if}"
%format then' = "\mathbf{then}"
%format else = "\mathbf{else}"
%format Par = "\constructor{Par}"
%format return = "\constructor{return}"
%format _>>=_ = _ "{" >>= "}\kern-\underscorespace" _
%format fail = "\constructor{fail}"
%format catch = "\constructor{catch}"
%format assert = "\constructor{assert}"
%format then = "\constructor{then}"
%format assert_then_ = assert _ then "\kern-\underscorespace" _
%link >>= _>>=_
%format Lens = "\constructor{Lens}"
%format Lens.put = Lens . put
%format Lens.get = Lens . get
%format Lens.PutGet = Lens . PutGet
%format Lens.GetPut = Lens . GetPut
%format put' = "\mathit{put}"
%format get' = "\mathit{get}"
%format ↔ = "\scalebox{.75}[1]{$\leftrightarrow$}"
%format _↔_ = _ "{" ↔ "}\kern-\underscorespace" _
%link ↔ _↔_
%format Iso = "\constructor{Iso}"
%format to-from-inverse = to - from - inverse
%format from-to-inverse = from - to - inverse
%format Iso.to = Iso . to
%format Iso.from = Iso . from
%format Iso.to-from-inverse = Iso . to-from-inverse
%format Iso.from-to-inverse = Iso . from-to-inverse
%format iso-lens = iso - lens
%format empty-iso = empty - iso
%format id-iso = id - iso
%format get-l-a↦b = get - l - a "{\mapsto}" b
%format put-r-b↦b' = put - r - b "{\mapsto}" b'
%format put-l-a-b'↦a' = put - l - a - b' "{\mapsto}" a'
%format get-r-b↦c = get - r - b "{\mapsto}" c
%format 2015 = "\mathrm{2015}"
%format <= = "{\leq}"
%format List = "\constructor{List}"
%format U = "\constructor{U}"
%format _⊕_ = _ "{" ⊕ "}\kern-\underscorespace" _
%link ⊕ _⊕_
%format _⊗_ = _ "{" ⊗ "}\kern-\underscorespace" _
%link ⊗ _⊗_
%format list = "\constructor{list}"
%format μ = "\mu"
%format Pattern = "\constructor{Pattern}"
%format var = "\constructor{var}"
%format k = "\constructor{k}"
%format one = "\constructor{one}"
%format left = "\constructor{left}"
%format right = "\constructor{right}"
%format prod = "\constructor{prod}"
%format unit = "\constructor{unit}"
%format interp-update = "\mathit{interp}" - "\mathit{update}"
%format Paths = "\constructor{Paths}"
%format BiGUL = "\constructor{BiGUL}"
%format failB = "\constructor{fail}"
%format skip = "\constructor{skip}"
%format replace = "\constructor{replace}"
%format update = "\constructor{update}"
%format rearr = "\constructor{rearr}"
%format dep = "\constructor{dep}"
%format caseS = "\constructor{caseS}"
%format caseV = "\constructor{caseV}"
%format align = "\constructor{align}"
%format BiGULs = "\constructor{BiGULs}"
%format CaseSBranch = "\constructor{CaseSBranch}"
%format CaseSBranch' = "\constructor{CaseSBranch^B}" 
%format CaseSBranchType = "\constructor{CaseSBranchType}"
%format CaseSBranchType' = "\constructor{CaseSBranchType^B}"
%format normal = "\constructor{normal}"
%format adaptive = "\constructor{adaptive}"
%format normal' = "\constructor{normal}"
%format adaptive' = "\constructor{adaptive}"
%format interp-caseS = "\mathit{interp}" - "\mathit{caseS}"
%format PutGet-with-adaptation = "\mathit{PutGet}" - "\mathit{with}" - adaptation
%format GetPut-with-adaptation = "\mathit{GetPut}" - "\mathit{with}" - adaptation
%format CaseVBranch = "\constructor{CaseVBranch}"
%format CaseVBranch' = "\constructor{CaseVBranch^B}" 
%format get-with-iso = "\mathit{get}" - "\mathit{with}" - iso
%format replace-lens = "\mathit{replace}" - lens
%format fail-lens = "\mathit{fail}" - lens
%format skip-lens = "\mathit{skip}" - lens
%format caseS-lens = "\mathit{caseS}" - lens
%format check-diversion = check - diversion
%format put-with-adaptation = put - "\mathit{with}" - adaptation
%format caseV-lens = "\mathit{caseV}" - lens
%format guarded-return = guarded - "\mathit{return}"
%format == = "\shorteq\kern.5pt\shorteq"
%format 0 = "\mathrm0"
%format source-condition = source - condition
%format matching-condition = matching - condition
%format align-lens = "\mathit{align}" - lens
%format Lenses = "\constructor{Lenses}"
%format Views = "\constructor{Views}"
%format Lenses' = Lenses "^\prime"
%format Views' = "\constructor{Views^B}"
%format update-lens = "\mathit{update}" - lens
%format U-dec = "\mathit{U}" - dec
%link ⟦ ⟦_⟧
%link ⟧ ⟦_⟧
%format Level = "\constructor{Level}"
%format VarPositions = "\constructor{VarPositions}"
%format PatResult = "\constructor{PatResult}"
%format pat-iso = pat - iso
%format Path = "\constructor{Path}"
%format eval-path = eval - path
%format view-rearr-iso = view - "\mathit{rearr}" - iso
%format Invertible = "\constructor{Invertible}"
%format bind-computes = bind - computes
%format relocate-employees = "\mathit{relocate}" - employees
%format relocate-employee = "\mathit{relocate}" - employee
%format adjust-salary-and-office = adjust - salary - "\mathit{and}" - office
%format relocate-to-UK = relocate - to - UK
%format relocate-to-US = relocate - to - US
%format toPounds = "{" $ "}{" ⇒ "}" £
%format £⇒$ = £ "{" ⇒ "}\kern-1.5pt" $

\newcommand{\shorteq}{%
  \resizebox{.4em}{1ex}{=}%
}
\newcommand{\ucomma}{\kern1pt, }
\newcommand{\uperiod}{\kern1pt. }

\newcommand{\shade}[1]{\colorbox[gray]{.9}{\smash{#1}\vphantom{\vrule height 1.1ex}}}
\newcommand{\lno}[1]{\text{\scriptsize #1\quad}}

\newlength{\underscorespace}
\setlength{\underscorespace}{.8pt}
\newlength{\binderspace}
\setlength{\binderspace}{.8pt}
\newlength{\binderskip}
\setlength{\binderskip}{.8pt}

\newcommand{\keyword}[1]{\mathbf{#1}}
\newcommand{\identifier}[1]{\mathit{#1}}
\newcommand{\constructor}[1]{\mathsf{#1}}

\usepackage[hidelinks]{hyperref}
\renewcommand{\chapterautorefname}{Chapter}
\renewcommand{\sectionautorefname}{Section}
\renewcommand{\subsectionautorefname}{Section}
\renewcommand{\figureautorefname}{Figure}
\def\equationautorefname~#1\null{(#1)\null}

\usepackage[color=yellow,textsize=footnotesize]{todonotes}

\usepackage{mathptmx}
\usepackage{amsmath}
\usepackage{amsthm}

\newtheorem*{definition}{Definition}
\newtheorem*{theorem}{Theorem}

\renewcommand{\emph}{\textsl}

\newcommand{\BiGUL}{\textsc{BiGUL}}
\newcommand{\Putlenses}{\textsc{Putlenses}}
\newcommand{\BiFluX}{\textsc{BiFluX}}
\newcommand{\Agda}{\textsc{Agda}}
\newcommand{\Haskell}{\textsc{Haskell}}
\newcommand{\BiYacc}{\textsc{BiYacc}}

\special{papersize=8.5in,11in}
\setlength{\pdfpageheight}{\paperheight}
\setlength{\pdfpagewidth}{\paperwidth}

\setlength{\mathindent}{.5\parindent}

\newcommand{\beforefigure}{\setlength{\mathindent}{0em}}

\hyphenation{Bi-Yacc}

\begin{document}

%\conferenceinfo{PEPM '16}{Month d--d, 20yy, City, ST, Country} 
%\copyrightyear{20yy} 
%\copyrightdata{978-1-nnnn-nnnn-n/yy/mm} 
%\doi{nnnnnnn.nnnnnnn}

% Uncomment one of the following two, if you are not going for the 
% traditional copyright transfer agreement.

%\exclusivelicense                % ACM gets exclusive license to publish, 
                                  % you retain copyright

%\permissiontopublish             % ACM gets nonexclusive license to publish
                                  % (paid open-access papers, 
                                  % short abstracts)

% \titlebanner{banner above paper title}        % These are ignored unless
% \preprintfooter{short description of paper}   % 'preprint' option specified.

\title{BiGUL: A Formally Verified Core Language for Putback-Based Bidirectional Programming}

\authorinfo{Hsiang-Shang Ko}
           {National Institute of Informatics, Japan}
           {hsiang-shang@@nii.ac.jp}
\authorinfo{Tao Zan \and Zhenjiang Hu}
           {SOKENDAI (The Graduate University for\\ Advanced Studies), Japan\\ National Institute of Informatics, Japan}
           {\{zantao,hu\}@@nii.ac.jp}

\maketitle

\todo[inline]{check grammar of conditional sentences}

\begin{abstract}

Putback-based bidirectional programming allows the programmer to write only one putback transformation, from which the unique corresponding forward transformation is derived for free.
The logic of a putback transformation is more sophisticated than that of a forward transformation and does not always give rise to well-behaved bidirectional programs; this calls for more robust language design to support development of well-behaved putback transformations.
In this paper, we design and implement a concise core language |BiGUL| for putback-based bidirectional programming to serve as a foundation for higher-level putback-based languages.
|BiGUL| is completely formally verified in the dependently typed programming language \Agda\ to guarantee that any putback transformation written in |BiGUL| is well-behaved.

\end{abstract}

%\category{CR-number}{subcategory}{third-level}

% general terms are not compulsory anymore, 
% you may leave them out
%\terms
%term1, term2

\keywords
putback-based bidirectional transformations,
dependent types,
formal verification

\section{Introduction}
\label{sec:introduction}
% Sketch:
%   1. Brief introduce BX, and then Start from Put-based BX
%   2. Talk the problem of Putlenses, BiFluX
%   3. Give our solution
%   4. Show the advantages

% Revise: emphasis on Formally verified,
%         also state Minimal but useful enough since it is a core for BiFluX.

% 1.
%   data replica
%   Internet of Things
%   information explosion
%   Maybe Remove this sentence.
%Due to the increase in technology, smart devices for different purpose are developed, thus come into an explosion in the number of devices and information.
%The duplication of information among different devices makes it hard to be kept consistent which is an important issue.
% BX for synchronization
\emph{Bidirectional transformations}~(BXs)~\cite{GRACE:09} provide a novel mechanism for maintaining consistency between two pieces of related information, one referred to as the source and the other as the view.
A bidirectional transformation consists of a pair of transformations ---
a forward |get| which extracts information from a source to construct an abstract view,
and a backward |put| which embeds information of a view back into a source, producing an updated source --- and this pair of transformations should be \emph{well-behaved}, i.e., they should satisfy two round-tripping laws (|PutGet| and |GetPut|; see \autoref{sec:BX}).
 % (source and view).
% From language approach: bx languages.
Since writing both transformations while guaranteeing their well-behavedness can take a lot of effort,
many bidirectional programming languages~\cite{Foster-lenses,Bohannon:06,MHNHT07,Hidaka:10,PachecoCunha:10,Barbosa:2010,Hofmann:2011,FoPP08,HoPW12} have been designed
to aid the user in writing bidirectional transformations,
with which the programmer only needs to write one program that can be interpreted either as |get| or |put|, and the two interpretations are guaranteed to be well-behaved.
%The correctness of this BX program is guaranteed by well-behavedness~\cite{Lenses}.

Recently there have been investigations into the feasibility of writing bidirectional programs by describing solely their putback directions~\cite{Hu-putback-validity, Pacheco-putlenses, Pacheco-BiFluX}.
The rationale behind this \emph{putback-based} approach is that the |put| component of a bidirectional transformation uniquely determines its |get| component~\cite[Lemma~2.2.5]{Foster:09}.
The combinator library \Putlenses~\cite{Pacheco-putlenses} follows this approach, providing a set of putback-based \emph{lens} combinators~\cite{Foster-lenses} for writing complex BX programs.
Based on \Putlenses\ and the functional XML update language \textsc{Flux}~\cite{Cheney-FLUX}, a bidirectional XML update language~\BiFluX~\cite{Pacheco-BiFluX} has also been designed and implemented.
Putback-based programming is more delicate than previous approaches which centre around writing |get|, not only because |put| --- being an update --- is inherently more complex than |get|, but also because not all updates give rise to well-behaved BXs.
This introduces a new challenge of ensuring the well-behavedness of a |put|, in the sense that a unique |get| can indeed be derived from the |put| to form a well-behaved BX.
\Putlenses\ and \BiFluX\ meet this challenge by inserting dynamic (runtime) checks around programmer-specified actions used in |put|, and the programmer is supposed to supply only actions that can pass these dynamic checks.
While both languages are carefully designed, the intricate nature of these dynamic checks can still make people cast doubt on their correctness.

%
%Most existing approaches to BXs are \emph{|get|-based}, in the sense that the system asks the programmer to write BX programs describing a |get|, % describing how to construct a view form a source
%and derives a backward |put| for free.\todo{``This is not entirely accurate, and gives the impression that techniques where the programmer writes both directions using constructs from which both directions can be derived (which is essentially the approach taken in this paper too) are little-investigated, which is not the case.''}\
%The major problem with the |get|-based approach is that there are usually more than one \emph{incomparable} |put|s for a |get|~\cite{BuCV08}\todo{check citation}, and thus the backward behaviour of a |get|-based program becomes ambiguous and unpredictable.
%%
%To remedy this situation, a dual \emph{putback-based} approach was proposed recently to support writing BXs in terms of |put| instead of |get|~\cite{Fischer-essence-of-BX}.
%The behaviour of |put| can thus be fully specified, and, unlike the |get|-based approach, the derived |get| is guaranteed to be unique and thus unambiguous.
%The update program user written using \textsc{BiFluX} in fact is the $put$ function, and the system derives $get$ for free.

% 2. Explain the problem
%    putlenses not fully checked ???
%    biflux not clearly checked, or hard to check ?
%    say: not fully, formally verified.

%
% gives opportunities to make BX really useful in real applications
%since it lets user to have full control over the update strategy.
%While the well-behavedness is not fully formal verified for both $putlenses$ and \textsc{BiFluX} whi%ch is the most important property for {BX}.
%As the forward transformation $get$ is derived for free
%and user 
%do not need to care how $get$ is derived 
%and not have to check the derived $get$ is correct or not manually, 
%it is necessary to fully guarantee the program user written is well-behaved by the BX language.
%
% concretely explain how it is in putlenses and biflux
% TODO: bad writing
%For example, in \textsc{BiFluX},
%the semantics of the basic update operator $replace$ is straight-forward
%and the well-beahvedness for $replace$ is not checked formally,
%and complex one like $if$ statement, the well-behavedness is checked dynamically,
%while all of them are not formally verified.
%Since \textsc{BiFluX} is designed for XML data,
%it also needs to handle complex regular expressions 
%and XML structured data which makes the program hard to check.

% Needs

In this paper, we aim to build a clean and solid foundation for higher-level putback-based languages (e.g., \BiFluX) by focusing on a small yet sufficiently powerful putback-based core language |BiGUL| (for \emph{Bidirectional Generic Update Language}), which is completely formally verified to guarantee that any |put| specified in the language is well-behaved.
In more detail:
\begin{itemize}
\item Firstly, |BiGUL| is a clean revision of the core of \BiFluX, and is designed to be closer to practical programming languages (than, e.g., \Putlenses).
It consists of a set of essential and orthogonal bidirectional statements, such as alignment of source and view lists, pattern matching--based independent source updates and invertible view computation, and case analyses on either source or view.
Although currently |BiGUL| is specified only in terms of its dependently typed abstract syntax (which is not intended to be programmer-friendly), |BiGUL| programs are pretty straightforward to write (or compile to), thanks to its native support of familiar programming constructs like pattern matching and case statements.
%In particular, source case analysis is revised to include the {\em adaptation} mechanism, so we can specify case analyses that can accept sources of more varied forms which previously cannot be dealt with.
% and make 
%it easy to verify well-behavedness.
% for specifying different update strategies.
% show more of BiGUL
% show more !

\item Secondly, the semantics of |BiGUL| is concretely defined as monadic programs, with all dynamic checks for guaranteeing well-behavedness explicitly appearing in the programs.
|BiGUL| is thus differentiated from the more abstract treatment by, e.g., Foster~et~al~\cite{Foster-lenses}, in which partiality is dealt with implicitly and various well-behavedness conditions are specified set-theoretically, creating a gap between the definitions and their implementations.
The concrete semantics of |BiGUL|, in contrast, can be faithfully ported to practical programming languages like \Haskell\ without having to fill the gap between abstract mathematics and concrete implementation.

\item Lastly, |BiGUL| and its well-behavedness has been completely formally modelled and verified using the dependently typed programming language \Agda~\cite{Norell-thesis,Nore08}.
(Full source code is available at \url{http://www.prg.nii.ac.jp/project/bigul}, which is checked by \Agda\ version 2.4.2.4 with standard library version~0.11.)
Since the semantics of |BiGUL| is defined as concrete programs and directly proved correct, we can highly confidently guarantee that these programs, when ported to other languages, are indeed well-behaved.
The main proof technique is also worth mentioning: we reify computation steps of monadic programs as easily manipulable data, so well-behavedness proofs --- which have to cover all possible execution traces by analysing program structure --- can be carried out just like doing ordinary functional programming.
\end{itemize}
% finally prove our language is well-behaved.
%
%We design \BiGUL\ to be minimal to make it easy for usage and property checking. 
%For example, there is no $if$ statement in \BiGUL\ as it can be expressed as a special case of $case$ statement.
%While the well-behavedness of \BiGUL\ is fully verified using \Agda\ language.

%   4. Show the advantages
%Our contribution can be summarised as follows:
%\begin{itemize}
%[leftmargin=0.3in]
%  \item we propose a minimalist putback-based bidirectional language \BiGUL\ as a core language which is fully certified to guarantee the correctness, and served as a base for other putback-based bidirectional languages.
%  \item A correct forward transformation $get$ is automatically constructed through \BiGUL\ programs. User concentrates on specifying $put$, our system guarantees the derived $get$ together with the $put$ is well-behaved.

For the rest of the paper, after explaining how to formalise and prove well-behavedness of BXs in \Agda\ in \autoref{sec:BX}, we present |BiGUL|, its proof-carrying bidirectional semantics, and some examples in \autoref{sec:BiGUL}, and conclude with a few remarks in \autoref{sec:concluding-remarks}.
This paper is typeset with a modified version of lhs2\TeX\ such that all global \Agda\ identifiers are hyperlinked to their definitions, which should be helpful to the readers who read this paper electronically.


\section{Formalisation of well-behaved BXs in \Agda}
\label{sec:BX}

Before presenting the |BiGUL| language in \autoref{sec:BiGUL}, we first give an account of our \Agda\ formalisation of the underlying notion of well-behaved partial bidirectional transformations, and also demonstrate by a simple example how well-behavedness proofs are constructed in our setting (\autoref{sec:sample-well-behavedness-proof}).
The key to the formalisation is the reification of computation steps of monadic programs as a type family~|_↦_| (\autoref{sec:reification}), which makes the construction of the proofs as easy as ordinary functional programming.

We follow Foster~et~al's classic \emph{lens} approach~\cite{Foster-lenses}, but for the |PutGet| and |GetPut| laws we adopt the variant used by Pacheco et~al~\cite{Pacheco-putlenses}.
A lens between a \emph{source} type~|S| and a \emph{view} type~|V| is defined in \Agda\ as the following record type:%
\footnote{Arguments wrapped in curly braces are \emph{implicit}, and are typeset in a less obtrusive style in this paper most of the time.
We do not need to supply the implicit arguments when applying a function if \Agda\ can infer what those arguments should be.}
\begin{code}
record ^^^Lens (S V : Set) : Set₁ where
  field
    ^^^put  : S → V → Par S
    ^^^get  : S → Par V
    ^^^PutGet  : (HIDE({s s' : S} {v : V} →))  (put s v ↦ s') → (get s' ↦ v)
    ^^^GetPut  : (HIDE({s : S} {v : V} →))     (get s ↦ v) → (put s v ↦ s)
\end{code}
That is, a lens is a pair of functions |put| and |get| satisfying the |PutGet| and |GetPut| laws.
The two laws are referred to as the \emph{well-behavedness} properties.
Precise definitions of |Par| and~|_↦_| will be presented later in \autoref{sec:reification}.
Here we offer an informal explanation first:
Both |put| and |get| are partial functions that may or may not successfully produce a result.
Execution of |put s v|, if successful, produces an updated version of the source~|s| by properly embedding all information of the view~|v| into~|s|, while |get s| extracts the view embedded in~|s|.
The |PutGet| law states that, for all |s|,~|s' : S| and |v : V|, if |put s v| successfully computes to~|s'|, then subsequently |get s'| should successfully compute to~|v|; this ensures that |put| does its job properly --- after |put| updates a source with a view, |get| can completely recover the view from the updated source.
On the other hand, the |GetPut| law states that, for all |s : S| and |v : V|, if |get s| successfully computes to~|v|, then subsequently |put s v| should successfully compute to~|s|; this ensures that |put| does nothing more than its designated job --- |put| performs no update on a source if the view it receives is directly extracted from the source by |get|.

Note that our definition of lenses is stronger than the one given by Foster~et~al~\cite{Foster-lenses}: our |PutGet| law says that successful computation of |put| guarantees success of the subsequent computation of |get|, whereas in Foster~et~al's definition there is no such guarantee; the |GetPut| law is similar.

\subsection{Reification of monadic combinators and computation steps}
\label{sec:reification}

The kind of partial transformations we need are actually total functions which wrap their results in the standard |^^^Maybe| datatype with two constructors |^^^just : A -> Maybe A| and |^^^nothing  : Maybe A| --- the result of a successful computation is wrapped in the |just| constructor, while failed computation is represented by |nothing|.
It is thus always possible to know whether a computation succeeds or not in a finite amount of time.
This notion of ``decidable partiality'' should be distinguished from the usual partiality as formalised in terms of complete partial orders and continuous functions between them:
when defining the semantics of |BiGUL| (more specifically, for view case analysis in \autoref{sec:view-case}), we need to write programs that produce some result or fail respectively when a sub-program fails or produces some result, but this is not possible if failed computation is represented as a least element of a complete partial order.
%(since such functions violate monotonicity, a necessary condition for computability).

We can choose to write our transformations directly as |Maybe|-programs, usually structured with monadic combinators~\cite{Moggi-monads,Wadler-essence} like ``return'' and ``bind''.
This approach does not work too satisfactorily, though (see the second half of \autoref{sec:sample-well-behavedness-proof} for an example), when we consider the kind of proofs we need to construct for these programs:
When proving |PutGet| and |GetPut|, the antecedent is a proof that a program --- consisting of sub-programs bound together by the combinators --- computes successfully, and we need to analyse it to obtain proofs saying that each of the sub-programs computes successfully, before these proofs are reassembled to show that a related composite program also computes successfully.
This task would be much easier if a proof of a successful composite computation is exactly a bunch of proofs that all its sub-computations succeed.
To achieve this, we can reify the monadic binding structure of |Maybe|-programs by deeply embedding the monadic combinators as the constructors of a datatype, so later we can define types of proofs of successful computation by induction on the binding structure.

We thus define a datatype |Par|, which is a deep embedding of the operators that we use to write partial transformations:%
\footnote{Identifiers with underscores can be used either in prefix or infix form.
For example, we can write either |_>>=_ mx f| or |mx >>= f|.}
\begin{code}
data ^^^Par : Set -> Set1 where
  ^^^return        : (HIDE({A : Set} ->))    A -> Par A
  ^^^_>>=_         : (HIDE({A B : Set} ->))  Par A -> (A -> Par B) -> Par B
  ^^^fail          : (HIDE({A : Set} ->))    Par A
\end{code}
There are also other deeply embedded operators for error handling (see \autoref{sec:view-case}) and assertion, which we omit here.
The meaning of the constructors of |Par| is formally given by the following interpreter to |Maybe|:%
\footnote{The |with| notation evaluates an expression (here |runPar mx|) and matches the result with an additional pattern (appearing on the right of~`$||$').}
\begin{code}
^^^runPar : (HIDE({A : Set} ->)) Par A -> Maybe A
runPar (return x) = just x
runPar (mx >>= f)  with runPar mx
runPar (mx >>= f)  | just x   = runPar (f x)
runPar (mx >>= f)  | nothing  = nothing
runPar fail = nothing
\end{code}
We say that |mx : Par A| computes to |x : A| exactly when |runPar mx ≡ just x|, and that |mx|~fails to compute exactly when |runPar mx ≡ nothing|.
The equality |runPar mx ≡ just x| is already a family of types of proofs of successful computation, but these equality proofs are not so straightforward to analyse and construct (see the second half of \autoref{sec:sample-well-behavedness-proof}).
We hence move on to define an equivalent family of types, whose proofs are easier to manipulate.

For every |mx : Par A| and |x : A|, we define a type |mx ↦ x| whose inhabitants explain how the sub-programs in |mx| compute successfully one by one, eventually producing~|x|:%
\footnote{We write dependent pair types $\mathrm{\Sigma}(x \in A)\;B\;x$ as |(x : A) × B x|.
The underscore in |(x : _) × (mx ↦ x) × (f x ↦ y)| leaves the type of~|x| for \Agda\ to infer.}
\begin{code}
^^^_↦_ : (HIDE({A : Set} →)) Par A → A → Set
(return x)  ↦ x'  =  x {-"\kern-2pt"-}≡{-"\kern-2pt"-} x'
(mx >>= f)  ↦ y   =  (x : _) × (mx ↦ x) × (f x ↦ y)
fail        ↦ x   =  ⊥
\end{code}
In prose: |return x| computes to~|x'| when |x|~is exactly~|x'|; |mx >>= f| computes to~|y| when there exists a value~|x| such that |mx|~computes to~|x| and |f x| computes to~|y|; and it is impossible for |fail| to compute successfully.
Here our use of the term ``computes to'' is justified, since we can prove that a term of type |mx ↦ x| can be constructed if and only if |runPar mx ≡ just x|, meaning that our reification of |runPar mx ≡ just x| as |mx ↦ x| is accurate.
We are thus entitled to use~|_↦_| in the statements of |PutGet| and |GetPut|.

\subsection{A sample well-behavedness proof}
\label{sec:sample-well-behavedness-proof}

To give the reader a taste of how well-behavedness proofs are constructed in our \Agda\ setting, let us try to define an operator~|_↔_| for lens composition:%
\footnote{A field of a record is extracted by applying the field's accessor function to the record.
For example, the accessor function |Lens.get| has type |(HIDE({S V : Set} ->)) Lens S V -> S -> Par V|.}
\begin{code}
^^^_↔_ : (HIDE({A B C : Set} →)) Lens A B → Lens B C → Lens A C
l ↔ r = record
  {    put  =  λ a c →  Lens.get l a >>= λ b →
                        Lens.put r b c >>= Lens.put l a
  SEP  get  = λ a → Lens.get l a >>= Lens.get r
  SEP  PutGet = ? SEP GetPut = ? }
\end{code}
The |put| and |get| functions for the composite lens |l ↔ r| can be pretty straightforwardly constructed from the |put| and |get| functions of the component lenses |l|~and~|r|.
As part of the definition of |_↔_|\ucomma we also need to prove that |PutGet| and |GetPut| hold for the pair of |put| and |get| we provide.
For |PutGet| we need to show that, for all |a|,~|a' : A|, and |c : C|, from a proof of |put a c MAPSTO a'| we can construct a proof of |get a' MAPSTO c|.
That is, we need to write a function that converts an inhabitant of the type |put a c MAPSTO a'| to one of the type |get a' MAPSTO c|.
\Agda\ automatically expands the two types by the definition of~|_↦_|\ucomma the first one to
\begin{code}
(b : B) × (Lens.get l a MAPSTO b) ×
  (b' : B) × (Lens.put r b c MAPSTO b') × (Lens.put l a b' MAPSTO a')
\end{code}
and the second one to
\begin{code}
(b : B) × (Lens.get l a' MAPSTO b) × (Lens.get r b MAPSTO c)
\end{code}
Thus all we need to do is convert a five-tuple to a three-tuple.
This is just ordinary functional programming: Applying |Lens.PutGet l| to the fifth component of type |Lens.put l a b' MAPSTO a'| we obtain an inhabitant of type |Lens.get l a' MAPSTO b'|, and applying |Lens.PutGet r| to the fourth component of type |Lens.put r b c MAPSTO b'| we obtain an inhabitant of type |Lens.get r b' MAPSTO c|.
Hence the proof term for |PutGet| is simply
\begin{equation}
|λ { (b , x , b' , y , z) → (b' , Lens.PutGet l z , Lens.PutGet r y) }|
\tag{$*$}
\label{lambda-expression}
\end{equation}
With a similar reasoning, we can also construct a proof term for |GetPut|:
\begin{code}
λ { (b , x , y) → (b , x , b , Lens.GetPut r y , Lens.GetPut l x) }
\end{code}
completing the definition of~|_↔_|\uperiod

For the rest of this paper we will omit the well-behavedness proofs, but a majority of these proofs --- while requiring more complicated case analyses --- are technically as simple as the above proofs for lens composition.

To emphasise the advantage of reasoning with |Par| and |_↦_|\ucomma let us consider how the well-behavedness proofs for~|_↔_| would be constructed if we write partial transformations simply as |Maybe|-programs and state well-behavedness in terms of \Agda's equality type~|_≡_|\uperiod
The |put| and |get| functions would be defined in exactly the same way except that the bind operator used is the |Maybe| version.
The |PutGet| law now reads
\begin{code}
^^^pg :  (HIDE({a a' : A} {c : C} →))
           (Lens.get l a >>= (λ b →
             Lens.put r b c >>= Lens.put l a) ≡ just a') →
           Lens.get l a' >>= Lens.get r ≡ just c
\end{code}
The logic of proving |pg| is exactly the same as before, but the machinery involved is more complex.
To analyse the antecedent equality proof, we can prove a lemma:
\begin{code}
bind-computes :
  (HIDE({A B : Set})) (mx : Maybe A) {f : A → Maybe B} {y : B} →
  mx >>= f ≡ just y →
  (x : A) × (mx ≡ just x) × (f x ≡ just y)
\end{code}
Applying the lemma twice analyses the antecedent into three smaller equality proofs, achieving the same effect as the pattern matching in the $\lambda$-expression~\autoref{lambda-expression}.
As for establishing the consequent, a probably easiest way is to rewrite |Lens.get l a'| to |just b'| with an application of |Lens.PutGet l| so the consequent type simplifies to |Lens.get r b' ≡ just c|, which can then be discharged by an application of |Lens.PutGet r|.
The whole proof is
\begin{code}
pg {a  = a  }  p with bind-computes (Lens.get l a) p
pg {c  = c  }  p | b , x , q with bind-computes (Lens.put r b c) q
pg             p | b , x , q | b' , y , z rewrite Lens.PutGet l z
{-"\hfill"-} = Lens.PutGet r y
\end{code}
which is apparently more heavyweight than~\autoref{lambda-expression}, even for this small lens composition example.
For more complicated well-behavedness proofs, the above approach quickly becomes unmanageable, whereas our approach scales nicely.

\todo[inline]{Hugo: In my understanding, you are advocating that a deep embedding of the proofs in Agda is more manageable than a shallow embedding? If so, this is standard nomenclature for PEPM, and deep embeddings are known for being more flexible.}

\section{BiGUL and its formally verified lens semantics}
\label{sec:BiGUL}

%Various bidirectional programming languages have been built based on the idea of lenses.
%Such languages would contain a set of primitive lenses and a set of combinators which combine smaller lenses into larger ones (like what |_↔_|~does).
%Usually these languages are designed to suggest that what they describe are |get| programs, but it has recently been argued that the opposite \emph{putback-based} approach should be adopted instead, due to the following theorem:\todo{previously appeared}
%\begin{theorem}[putback is the essence]
%Define, for all |mx|,~|my : Par A|, a type |mx ^^^≈ my| to mean that |mx MAPSTO z| if and only if |my MAPSTO z| for any |z : A|.
%Then for any two lenses |l|,~|r : Lens S V| whose |put| components are extensionally equal, i.e.,
%|Lens.put l s v ≈ Lens.put r s v| for all |s : S| and |v : V|, their |get| components must also be extensionally equal, i.e., |Lens.get l s ≈ Lens.get r s| for all |s : S|.
%\end{theorem}
%The theorem tells us that, when designing a bidirectional combinator, if we can justify the design of its |put| component, then its |get| component needs no further justification, since there is no other |get| that can be paired with the same |put|.
%This is in sharp contrast with the |get|-based approach, with which the designer needs to choose for each |get| one or a few possible |put| strategies to be paired with the |get|.
%This complicates the design as there would be multiple combinators with the same |get| semantics, distinguished only by differences in their |put| semantics --- this is rather counterintuitive as the programmer is supposed to think in terms of |get| instead of |put| when programming in |get|-based languages.
%Also it may well happen that, after the programmer describes a |get|, the corresponding |put| synthesised by the language is not what the programmer expects.
%All these problems simply do not exist for the putback-based approach.

\begin{figure}[t]
\beforefigure
\begin{code}
data ^^^BiGUL : U → U → Set₁ where
  ^^^replace  :  (HIDE({S : U} →))    BiGUL S S
  ^^^failB    :  (HIDE({S V : U} →))  BiGUL S V
  ^^^skip     :  (HIDE({S : U} →))    BiGUL S one
  ^^^caseS    :  (HIDE({S V : U} →))  List (CaseSBranch'  S V) → BiGUL S V
  ^^^caseV    :  (HIDE({S V : U} →))  List (CaseVBranch'  S V) → BiGUL S V
  ^^^align    :  (HIDE({S V : U} →))  (source-condition : ⟦ S ⟧ → Par Bool)
                                      (match : ⟦ S ⟧ → ⟦ V ⟧ → Par Bool)
                                      (b : BiGUL S V)
                                      (create : ⟦ V ⟧ → Par ⟦ S ⟧)
                                      (conceal : ⟦ S ⟧ → Par (Maybe ⟦ S ⟧)) →
                                      BiGUL (list S) (list V)
  ^^^update   :  (HIDE({S : U} →))    (p : Pattern S) (bs : BiGULs p) →
                                      BiGUL S (Views p bs)
  ^^^rearr    :  (HIDE({S V V' : U} →))  (p : Pattern V) (q : Pattern V') →
                                         Paths p q → BiGUL S V' → BiGUL S V
\end{code}
\caption{Top-level definition of |BiGUL| (simplified)}
\label{fig:BiGUL}
\end{figure}

Having explained the infrastructure of our \Agda\ formalisation, we now move on to the |BiGUL| language itself.
Following the putback-based approach, our language |BiGUL| describes |put| functions, i.e., updates to a source using a view.
\autoref{fig:BiGUL} shows (a simplified version of) the main definition of |BiGUL| (which refers to some auxiliary definitions that will appear in later sub-sections).
In general, a |BiGUL| update proceeds by decomposing the source into several parts to be updated independently~(|update|; \autoref{sec:source-update}), and accordingly rearranging the view to match with these parts~(|rearr|; \autoref{sec:view-rearrangement}), until the source and the view are reduced to the extent that basic operations can be applied~(|replace|, |failB|, and |skip|; \autoref{sec:basic-lenses}).
When both the source and the view are lists, there is a powerful alignment operation for synchronising the two lists~(|align|; \autoref{sec:list-alignment}).
And we can do case analyses on either the source~(|caseS|; \autoref{sec:source-case}) or the view~(|caseV|; \autoref{sec:view-case}) for describing more complex update strategies.

\begin{figure*}
\beforefigure
\begin{code}
{-"\lno1"-}  align  (λ { (_ , _ , year , _ , instock) → year == 2015 && instock })
{-"\lno2"-}         (λ { (stitle , _) (vtitle , _) → stitle == vtitle })
{-"\lno3"-}         (rearr  (prod var var)  (prod  var (prod    unit (prod   unit (prod  var          unit))))
{-"\lno4"-}                                 (      inj₁ refl ,  tt ,         tt ,        inj₂ refl ,  tt)
                    {-""-}  -- the two lines above are a deeply embedded representation of the function |λ { (title , price) -> (title , tt , tt , price , tt) }|
{-"\lno5"-}         (update  (prod  var (prod        var (prod     var (prod     var              var))))
{-"\lno6"-}                  ((_ ,  replace) , (_ ,  skip) , (_ ,  skip) , (_ ,  replace) , (_ ,  skip))))
{-"\lno7"-}         (λ _ -> return ("" , "(to be updated)" , 2015 , 0 , true))
{-"\lno8"-}         (λ { (title , author , year , price , instock) -> return (just (title , author , year , price , false)) })
\end{code}
\caption{Updating the prices of the books published in 2015}
\label{fig:bookstore}
\end{figure*}

\autoref{fig:bookstore} gives a preview of a |BiGUL| program.
(A more complex example, in particular involving |caseS| and |caseV|, is given in \autoref{sec:transatlantic-corporation}.)
The source is a list of book data of type |String × String × ℕ × ℕ × Bool| representing a book's title, author, publication year, price, and whether it is in stock, and we intend to update the price of those books published in 2015 which are still in stock.
We thus prepare a list of pairs of book title and price of type |String × ℕ| as the view.
The |BiGUL| program then describes how to update the source with the view:
we focus on those books in the source list which are published in 2015 and still in stock~(line~1) and align them with the view list by title~(line~2);
for every matched pair of source and view books, we rearrange the view~(lines 3--4) to match with the shape of the source, decompose the source by pattern matching (line~5), and replace the title and price of the source with those from the view~(line~6);
for every unmatched view book, we create a source book with default author information and publication year 2015 and mark it as in stock~(line~7), and its title and price will later be updated by the action for matched pairs (i.e., lines 3--6);
finally, for every unmatched source book, we keep it in the source list but mark it as out of stock~(line~8).

The types of the |BiGUL| constructors are indexed by the source and view types of the operations represented by the constructors, and ultimately we can write a well-typed interpreter to |Lens|:
\begin{code}
^^^interp : (HIDE({S V : U} ->)) BiGUL S V -> Lens ⟦ S ⟧ ⟦ V ⟧
\end{code}
(|U|~and~|⟦_⟧| are defined in \autoref{sec:source-update}.)
Since well-behavedness is built into the definition of |Lens|, being able to construct |interp| means that we have not only translated |BiGUL| to |put| and |get| functions, but also proved formally that |PutGet| and |GetPut| are satisfied, in the way explained in \autoref{sec:sample-well-behavedness-proof}.
For each operation we will start from a natural putback semantics, then add necessary checks so a corresponding |get| can exist, arriving at the lens semantics for the operation (including well-behavedness proofs), and finally deeply embed the lens as a constructor of |BiGUL|.
(An exception is |update|, the development for which deviates slightly from the general process to avoid some typing problem.)
We will go through the whole process for the three basic operations in \autoref{sec:basic-lenses} and source case analysis in \autoref{sec:source-case}.
For the rest of the operations we will only give a sketch for brevity.


\subsection{Replacing, failure, and skipping}
\label{sec:basic-lenses}

The three most basic operations are |replace|, |failB|, and |skip|.
The |replace| operation is applicable when the source and view are of the same type.
Its |put| semantics is discarding the original source and returning the view as the updated source, while its |get| semantics is just the identity.
This bidirectional semantics can in fact be derived from a \emph{partial isomorphism}:
We define partial isomorphisms by
\begin{code}
record ^^^Iso (A B : Set) : Set₁ where
{-"~"-}  field
           ^^^to    : A → Par B
           ^^^from  : B → Par A
           ^^^to-from-inverse  : (HIDE({x : A} {y : B} →)) (to x ↦ y) → (from y ↦ x)
           ^^^from-to-inverse  : (HIDE({x : A} {y : B} →)) (from y ↦ x) → (to x ↦ y)
\end{code}
of which |Lens| is a generalisation, as we can easily convert |Iso A B| to |Lens A B|:
\begin{code}
^^^iso-lens : (HIDE({A B : Set} →)) Iso A B → Lens A B
iso-lens iso = record  {    put     = λ _ b → Iso.from iso b
                       SEP  get     = Iso.to iso
                       SEP  PutGet  = Iso.from-to-inverse iso
                       SEP  GetPut  = Iso.to-from-inverse iso }
\end{code}
The lens semantics of |replace| can then be derived from the identity isomorphism:
\begin{code}
^^^id-iso : (HIDE({S : Set} ->)) Iso S S
id-iso = record  { to = return SEP from = return SEP ABBRV }

interp replace = iso-lens id-iso
\end{code}
We can also derive the lens semantics for |failB| --- whose |get| and |put| simply fail to compute for any input --- from the empty isomorphism in the same way:
\begin{code}
^^^empty-iso : (HIDE({S V : Set} ->)) Iso S V
empty-iso = record  {    to    = \ _ -> fail
                    SEP  from  = \ _ -> fail SEP ABBRV }

interp failB = iso-lens empty-iso
\end{code}
On the other hand, the |put| semantics of the |skip| operation is discarding the view and returning the original source, and its lens semantics is not an instance of |iso-lens|:
\begin{code}
^^^skip-lens : (HIDE({S : Set} ->)) Lens S ⊤
skip-lens = record  {    put  = \ s _ -> return s
                    SEP  get  = \ _ -> return tt SEP ABBRV }

interp skip = skip-lens
\end{code}
Note that the view type is specified as the unit type~|^^^⊤| (whose only inhabitant is~|^^^tt|) --- had the view type been more complex, there would have been no way for |get| to decide what to return.
In general, |put| must embed all view information into the updated source so |get| can recover the entire view from the updated source, establishing |PutGet|.
|⊤|~is a type with no information, and hence |Lens.put skip-lens| can safely ignore its view argument of type~|⊤|.

\begin{code}
\end{code}

\subsection{Case analysis on source}
\label{sec:source-case}

\begin{figure*}
\beforefigure
\begin{code}
^^^caseS-lens : (S V : Set) -> List (CaseSBranch S V) -> Lens S V
caseS-lens S V bs = record  {    put  = \ s v -> put-with-adaptation bs [] s v (λ s' → put-with-adaptation bs [] s' v (\ _ -> fail))
                            SEP  get  = get' bs SEP ABBRV }
  where
    ^^^branch : (HIDE'({l : Level})) {X : Set (HIDE'(l))} → (Lens S V → X) → ((S → Par S) → X) → CaseSBranchType → X
    branch f g (normal    l  ) = f  l
    branch f g (adaptive  u  ) = g  u

    ^^^check-diversion : List (CaseSBranch S V) → S → Par ⊤
    check-diversion  []              s = return tt
    check-diversion  ((p , _) ∷ bs)  s = p s >>= λ matched → if matched then' fail else check-diversion bs s

    ^^^put-with-adaptation : List (CaseSBranch S V) → List (CaseSBranch S V) → S → V → (S → Par S) → Par S
    put-with-adaptation []              bs' s v cont = fail
    put-with-adaptation ((p , b) ∷ bs)  bs' s v cont =
      p s >>= λ matched → if matched   then'  branch  (λ l →  Lens.put l s v >>= λ s' → p s' >>= λ matched' →
                                                              if matched' then' check-diversion bs' s' >>= (\ _ → return s') else fail)
                                                      (λ u → u s >>= cont) b
                                       else   put-with-adaptation bs ((p , b) ∷ bs') s v cont

    get' : List (CaseSBranch S V) → S → Par V
    get' []              s = fail
    get' ((p , b) ∷ bs)  s = p s >>= λ matched → if matched then' branch (λ l → Lens.get l s) (\ _ -> fail) b else get' bs s
    
    ABBRV
\end{code}
\caption{Definitions of |put| and |get| for source case analysis}
\label{fig:caseS}
\end{figure*}

The next construct |caseS| is for doing case analysis on the source.
The basic idea of |caseS|'s semantics is to choose from a list of branches the first one that matches the source, and then execute the |BiGUL| statement of that branch.
It is, however, difficult to find a |get| to pair with this |put| semantics: the obvious |get| which selects the first branch that matches its source argument does not work, because, when we consider |PutGet|, it may well happen that the updated source produced by |put| matches a different branch from the one matching the original source, and hence in general we would have to guarantee |PutGet| for |put| and |get| coming respectively from any two branches, which is unmanageable.
A more manageable solution is to insert a dynamic check at the end of |put| to ensure that the updated source must match the same branch as the original source, which is the solution adopted by Foster~et~al's concrete conditional lens~\cite[Section~6.1]{Foster-lenses}.

The above solution, however, is sometimes too restrictive in practice.
For example, in our putback-based language \BiYacc~\cite{Zhu-BiYacc} for generating well-behaved pairs of parsers and ``reflective'' printers, a program specifies how a printer \emph{updates} a concrete syntax tree (the source) to become consistent with an abstract syntax tree (the view).
The printer tries to retain the structure of a source wherever possible, but when the structure of the source is radically different from that of the view, we need to change the source completely, and this change of source structure would fail the aforementioned dynamic check that prevents branch switching.
We therefore need a more flexible case analysis on source.
(See \autoref{sec:transatlantic-corporation} for a simpler scenario where we also need this increased flexibility.)

\subsubsection{Enriching source case analysis with adaptive branches}
Our solution is to add a special kind of branches called \emph{adaptive} branches to source case analysis, in addition to \emph{normal} branches.
The semantics of the two kinds of branches are defined by the following datatype:
\begin{code}
data ^^^CaseSBranchType (S V : Set) : Set1 where
  ^^^normal    : Lens S V      -> CaseSBranchType S V
  ^^^adaptive  : (S -> Par S)  -> CaseSBranchType S V
\end{code}
The semantics of a |normal| branch is just a lens, while for an |adaptive| branch it is a source transformation.
With each branch we also need to associate a decidable predicate on the source type specifying when a source matches the branch, so the complete definition of the type of branches is
\begin{code}
^^^CaseSBranch : Set -> Set -> Set1
CaseSBranch S V = (S -> Par Bool) × CaseSBranchType S V
\end{code}
The lens semantics of |caseS| is then given by
\begin{code}
caseS-lens : (S V : Set) -> List (CaseSBranch S V) -> Lens S V
\end{code}
which is built from a list of branches.

\subsubsection{The \textit{put} and \textit{get} semantics}

The |put| and |get| components of |caseS-lens| are shown in \autoref{fig:caseS}.
For the |put| semantics, we might think of a source case analysis as still consisting mainly of normal branches, but we can now specify additional cases such that when a source matches none of the normal branches, the source can still be accepted, which is then transformed --- or adapted --- to one that will match a normal branch.
More explicitly, the case analysis may be run twice: if a normal branch is chosen the first time, then the case analysis succeeds and we execute the branch; otherwise, we adapt the source and rerun the case analysis, and must match a normal branch this time.
A single round of the case analysis is performed by the function |put-with-adaptation|, which takes a continuation of type |S -> Par S| that, when an adaptive branch is matched, is invoked on the adapted source.
The |put| semantics, which consists of up to two rounds of the case analysis, is |put-with-adaptation| with a second instance of |put-with-adaptation| as the continuation, and the continuation for the latter is the computation that always fails --- that is, the second round cannot fall into an adaptive branch.
In either round, if a normal branch is matched and executed, we must ensure that the updated source satisfies the predicate of the branch and also none of the predicates of the previous branches, so a subsequent execution of |get| on the updated source --- which chooses a branch in the same way as |put| --- would go through the same branch as |put|, establishing |PutGet|.
In |put-with-adaptation|, there is an accumulating parameter~|bs'| keeping hold of the mismatched branches; if a normal branch is matched and executed, the function |check-diversion| would be invoked to ensure that the updated source matches none of the branches in~|bs'|.

The |get| semantics, on the other hand, simply tries to match the source with each of the branches and execute the first matched branch if the branch is normal, or fail if the branch is adaptive.
A quick reasoning for the behaviour about adaptive branches is as follows:
An abstract reading of the well-behavedness laws says that the domain of |get| must coincide with the range of updated sources produced by |put|.
Since a successfully updated source necessarily comes out from, and will again match, a normal branch, |get| must fail for any source matching an adaptive branch.
Of course, we still need to prove the well-behavedness properties to actually say that the |put| and |get| functions are defined correctly.

\subsubsection{Sketch of the well-behavedness proofs}

As explained in \autoref{sec:sample-well-behavedness-proof}, the well-behavedness proofs for |caseS-lens| proceed by analysing all possible ways for |put| or |get| to compute successfully, and then showing how their successful computation guarantees successful computation of the corresponding |get| or |put|.
The only slight complication about |caseS-lens| is that |put-with-adaptation| has an accumulating parameter, which implies that we need to prove generalised versions of the well-behavedness statements.
For |PutGet|, the main statement we prove is
\begin{code}
PutGet-with-adaptation :
  (bs bs' : List (CaseSBranch S V))
  {s s' : S} {v : V} {cont : S → Par S} →
  ({s : S} →  (cont s ↦ s') → (get' (revcat bs' bs) s' ↦ v)) →
  (put-with-adaptation bs bs' s v cont ↦ s') →
  (get' (revcat bs' bs) s' ↦ v)
\end{code}
where |revcat| is defined by
\begin{code}
^^^revcat : (HIDE'({l : Level})) {A : Set (HIDE'(l))} → List A → List A → List A
revcat []        ys = ys
revcat (x ∷ xs)  ys = revcat xs (x ∷ ys)
\end{code}
which is the usual way of implementing linear-time list reversal, and coincides with how branches are moved from |bs| to |bs'| in |put-with-adaptation|.
%Here is one way of thinking about this generalised statement:
%During execution of |put|, the list of branches is divided into |bs| and~|bs'|, respectively containing those branches yet to be matched and those mismatched with the source, and there is an invariant that the original list of branches is always |revcat bs' bs|.
%|PutGet-with-adaptation| then guarantees that |get| using the original list of branches will succeed if |put-with-adaptation| succeeds and the continuation satisfies a |PutGet|-like property (so the responsibility of establishing |PutGet| can be delegated to this proof if the continuation is invoked), and can be proved by induction on~|bs|.
|PutGet| is then proved by setting |bs'| to~|[]| and applying |PutGet-with-adaptation| to itself in a way similar to how we define |put| in terms of |put-with-adaptation|.
For |GetPut|, the generalisation we need is relatively simpler:
\begin{code}
GetPut-with-adaptation :
  (bs bs' : List (CaseSBranch S V))
  {cont : S → Par S} {s : S} {v : V} →
  (check-diversion bs' s ↦ tt) →
  (get' bs s ↦ v) → (put-with-adaptation bs bs' s v cont ↦ s)
\end{code}
%It states that if |get| using the list~|bs| of branches succeeds, then |put-with-adaptation| will also succeed as long as all branches in~|bs'| indeed are mismatched branches, i.e., |check-diversion| is successful.
|GetPut| is then proved by setting |bs'| to |[]|, in which case the premise about |check-diversion| becomes trivially true.

\subsubsection{Fitting the lens to the interpreter}

The final step is to deeply embed |caseS-lens| as a constructor of the |BiGUL| datatype.
We start with defining a variant of |CaseSBranchType| whose source and view type parameters are codes from the universe~|U| (\autoref{sec:source-update}) and whose |normal'| constructor takes a |BiGUL| statement instead of a |Lens|:
\begin{code}
data ^^^CaseSBranchType' (S V : U) : Set1 where
  ^^^normal'    : BiGUL S V             → CaseSBranchType' S V
  ^^^adaptive'  : (⟦ S ⟧ -> Par ⟦ S ⟧)  → CaseSBranchType' S V
\end{code}
and also a variant of |CaseSBranch|:
\begin{code}
^^^CaseSBranch' : U -> U -> Set1
CaseSBranch' S V =
  (⟦ S ⟧ -> Par Bool) × CaseSBranchType' S V
\end{code}
As shown in \autoref{fig:BiGUL}, the |caseS| constructor of |BiGUL| takes a list of |CaseSBranch'|s.
Its interpretation is, unsurprisingly, |caseS-lens|:
\begin{code}
mutual
  ABBRV
  interp (caseS {S} {V} bs) =
    caseS-lens ⟦ S ⟧ ⟦ V ⟧ (interp-caseS bs)
  ABBRV
  
  ^^^interp-caseS : (HIDE({S V : U} →))  List (CaseSBranch'   S      V    ) →
                                         List (CaseSBranch ⟦  S ⟧ ⟦  V ⟧  )
  interp-caseS []                                    = []
  interp-caseS (  (p , normal'  b           ) ∷ bs)  =
                  (p , normal   (interp b)  ) ∷ interp-caseS bs
  interp-caseS (  (p , adaptive'  u         ) ∷ bs)  =
                  (p , adaptive   u         ) ∷ interp-caseS bs
\end{code}
The helper function |interp-caseS| simply maps |interp| into the normal branches.
To help \Agda\ see that |interp| is terminating, though, we need to define |interp-caseS| explicitly (instead of using the standard |map| function on lists) and put the definition along with |interp| in a |mutual| block.

\subsection{Case analysis on view}
\label{sec:view-case}

For |caseV|, basically we still try to match the view with a list of branches, and execute the first branch that matched.
We must be more careful when it comes to the view, though:
As mentioned at the end of \autoref{sec:basic-lenses}, we must ensure that view information is completely embedded into the source, and that requires conscious management of view information.
When the view matches a branch, the amount of view information that remains to be embedded into the source may actually be reduced.
For example, if there is a branch that matches an integer view which equals zero, then for this branch we may even perform no update on the source as long as we can guarantee that the subsequent |get| will choose this branch, in which case |get| can simply return zero --- the information that the view is zero is embedded implicitly in the control flow rather than explicitly in the updated source.
Thus, in the definition of |CaseVBranch| below, we use a partial isomorphism on the view type which converts it to a possibly less informative type, instead of just a decidable predicate as in |CaseSBranch|:
\begin{code}
^^^CaseVBranch : Set -> Set -> Set₁
CaseVBranch S V = (V' : Set) × Lens S V' × Iso V' V
\end{code}
For the comparison-with-zero example above, we would set~|V'| to~|⊤|, indicating that there is no more information to be embedded into the source after a successful matching, and set the |to| and |from| components of the isomorphism to |\ _ -> return 0| and |\ n -> if n == 0 then' return tt else fail|.
A branch is considered matched with a view when the |from| conversion of the associated isomorphism succeeds on the view.
(The deep embedding as |BiGUL|'s |caseV| constructor restricts the isomorphisms to those induced by pattern matching (to be introduced in \autoref{sec:source-update}) --- a deeply embedded branch of type |^^^CaseVBranch' S V| is a pair of a pattern and a |BiGUL| statement that updates a source of type~|S| with the result of decomposing a view of type~|V| using the pattern.)

\begin{figure}[t]
\begin{code}
-- Assume |S : Set| and |V : Set|

get-with-iso : {V' : Set} → Lens S V' → Iso V' V → S → Par V
get-with-iso l iso s = Lens.get l s >>= Iso.to iso

put' : (bs : List (CaseVBranch S V)) → S → V → Par S
put' []                    s v = fail
put' ((_ , l , iso) ∷ bs)  s v =
  catch  (Iso.from iso v) (Lens.put l s)
         (  put' bs s v >>= λ s' →
              catch  (get-with-iso l iso s') (\ _ -> fail)
                     (return s'))

get' : (bs : List (CaseVBranch S V)) → S → Par V
get' []                    s = fail
get' ((_ , l , iso) ∷ bs)  s =
  catch  (get-with-iso l iso s) return
         (  get' bs s >>= λ v → catch  (Iso.from iso v) (\ _ -> fail)
                                       (return v))
\end{code}
\caption{Definitions of |put| and |get| for view case analysis}
\label{fig:caseV}
\end{figure}

The lens semantics we assign to |caseV| has type
\begin{code}
caseV-lens :  (S V : Set) -> List (CaseVBranch S V) -> Lens S V
\end{code}
The definitions of its |put| and |get| are shown in \autoref{fig:caseV}.
They require |Par| to be extended with one more constructor\todo{Hugo: perhaps make a reference to monadic exceptions or MonadPlus to explain the intuition of catch}
\begin{code}
^^^catch : (HIDE({A B : Set} →)) Par A → (A → Par B) → Par B → Par B
\end{code}
interpreted by
\begin{code}
runPar (catch mx f my) with runPar mx
runPar (catch mx f my) | just x   = runPar (f x)
runPar (catch mx f my) | nothing  = runPar my
\end{code}
That is, |catch mx f my| behaves like |mx >>= f| when |mx| computes successfully, or becomes |my| when |mx|~fails.
For |put|, then, we can choose to execute the current branch or try the rest of the branches based on whether |Iso.from iso| computes successfully.

For |get|, there is no obvious way to decide which branch to go --- unlike |put|, we are not given a view to match with the branches.
Here we try to execute each of the branches until we reach one that computes a view successfully, which is the approach taken by Pacheco~et~al~\cite{Pacheco-putlenses}.
This makes guaranteeing |PutGet| difficult, as we cannot easily guarantee that executing |put| followed by |get| would go through the same branch.
Our compromised solution is to make |put| do an expensive check after a branch is matched and executed, requiring that the |get| for each of the previous branches must fail.
To pass this check, the ranges of the |put| semantics of the branches (which coincide with the domains of |get|) have to be disjoint in order for |put| to compute successfully.

\subsection{List alignment}
\label{sec:list-alignment}

For the semantics of |align|, which is the key operation used in \autoref{fig:bookstore} (and in the \BiFluX\ language~\cite{Pacheco-BiFluX}), we employ a particular parametrised strategy for updating a list of sources of type~|S| with a list of views of type~|V|, which tries to align/match the sources with the views, synchronise the matched pairs of sources and views by an inner lens, and process the unmatched sources and views in some programmer-specified ways.
The purpose of |align| is similar to Barbosa~et~al's matching lenses~\cite{Barbosa:2010}, but we do not intend |align| to be as expressive as matching lenses yet --- a few of the design choices made below will appear somewhat arbitrary and inflexible, but the the semantics is already interesting enough such that formal verification is desirable.

The lens semantics of |align| has type
\begin{code}
align-lens : (HIDE({S V : Set} ->))
  (source-condition : S → Par Bool)
  (matching-condition : S → V → Par Bool)
  (l : Lens S V)
  (create : V → Par S)
  (conceal : S → Par (Maybe S)) →
  Lens (List S) (List V)
\end{code}
Its |put| semantics is as follows:
First, from the source list we can select a sub-list that we intend to align with the views by specifying a decidable predicate 
|source-condition : S -> Par Bool|.
The sub-list of sources satisfying the |source-condition| are then aligned with the list of views by finding pairs of source and view satisfying a given binary predicate |matching-condition : S -> V -> Par Bool|.
The matching order is fixed in our strategy: for each view we find the first matching source that has not been matched with a previous view.
On each matched pair of source and view we execute an inner lens |l : Lens S V| to update the source with the view.
For unmatched views, we should create corresponding sources in the source list (so a subsequent |get| can produce these views).
This is done by first applying a programmer-specified function |create : V → Par S| to the view to create a temporary source, and then updating this source with the view by~|l|.
Finally, for unmatched sources, we can choose to simply delete them or transform them such that they do not satisfy the |source-condition|.
This is specified by a function |conceal : S → Par (Maybe S)|: an unmatched source is deleted if applying |conceal| to it produces |nothing|, or transformed if |conceal| produces a new source wrapped in |just|.
These transformed sources can be placed anywhere in the source list; currently we simply place all of them at the front of the list.
Finally, this list of updated sources is merged with those not satisfying |source-condition| in the beginning.

The corresponding |get| extracts from the input source list all the sources satisfying |source-condition| and then uses |Lens.get l| to get a list of views from these sources.
To guarantee well-behavedness, a few checks about |l|~and |conceal| need to be inserted:
After a pair of source and view are synchronised by~|l|, the updated pair should again satisfy |matching-condition|, and also the updated source should still satisfy |source-condition|.
And a source produced by |conceal| must not satisfy |source-condition|.

\subsection{Source update}
\label{sec:source-update}

An indispensable operation is decomposing a source and updating its parts, represented by the |update| constructor of |BiGUL|.
We follow the approach taken by \textsc{Flux}~\cite{Cheney-FLUX}, which requires that updates must be independent --- that is, the same part of a source cannot be updated more than once --- to avoid the problem of sequencing updates.
It turns out that a pattern matching notation suits this task perfectly: pattern matching is arguably the most intuitive way to decompose a source, and we can specify apparently independent updates by writing them at the variable positions in a pattern.
For example, in \autoref{fig:bookstore} we use |update| to match a source book of type |String × String × ℕ × ℕ × Bool| with a five-variable pattern and update its five components respectively by |replace|, |skip|, |skip|, |replace|, and |skip|.
The idea itself is straightforward; what follows, despite its slight complexity (especially if the reader is not familiar with dependently typed programming), is merely our datatype-generic and strongly typed implementation of the idea in \Agda~\cite{Gibbons-DGP,Altenkirch-GP-within-DTP}.
The definitions for |update| will be laid out in three levels:
we first define the class of datatypes that we wish to process with |BiGUL|; for every datatype, we define the patterns applicable to the inhabitants of the datatype; finally, every pattern induces a ``container'' type, whose inhabitants can be used to store inner |BiGUL| statements at the variable positions of the pattern.

Since \Agda\ is fully dependently typed, we can naturally define patterns in a type-directed way such that nonsensical patterns are syntactically ruled out.
This task starts from the construction of a \emph{universe}~|U|, i.e., a datatype whose inhabitants are interpreted as types.
|U| has appeared in the definition of |BiGUL| in \autoref{fig:BiGUL}, where it is used as the type of the indices representing the source and view types of the operations.
The actual \Agda\ implementation of |BiGUL| uses a universe capable of expressing mutually inductive datatypes, but, to make the presentation easier to follow, the universe~|^^^U| that we define below --- along with its interpreter~|^^^⟦_⟧| --- is only a much simplified version:
\begin{code}
data U : Set₁ where             ⟦_⟧ : U → Set
  ^^^k     : Set → U    COLSEP  ⟦ k A     ⟧ = A
  ^^^one   : U                  ⟦ one     ⟧ = ⊤
  ^^^_⊕_   : U → U → U          ⟦ F ⊕ G   ⟧ = ⟦ F ⟧  ⊎  ⟦ G ⟧
  ^^^_⊗_   : U → U → U          ⟦ F ⊗ G   ⟧ = ⟦ F ⟧  ×  ⟦ G ⟧
  ^^^list  : U → U              ⟦ list F  ⟧ = List ⟦ F ⟧
\end{code}
The types that we can encode within~|U| are those expressible in terms of atomic types, the unit type~|⊤|, sum types~|_⊎_| (whose two constructors are |^^^inj₁| and |^^^inj₂|), product types~|_×_|\ucomma and |List|; applying |⟦_⟧| to an inhabitant of the universe decodes it to the type it represents.

We can now define a family of types |Pattern : U -> Set| such that, for any |D : U|, the type |Pattern D| contains exactly those patterns sensible for |⟦ D ⟧|.
Below is a simplified version covering the patterns used in this paper:
\begin{code}
data ^^^Pattern : U → Set₁ where
  ^^^var    : (HIDE({D : U} →)) Pattern D
  ^^^unit   : Pattern one
  ^^^left   : (HIDE({D E : U} →)) Pattern D → Pattern (D ⊕ E)
  ^^^right  : (HIDE({D E : U} →)) Pattern E → Pattern (D ⊕ E)
  ^^^prod   : (HIDE({D E : U} →)) Pattern D → Pattern E → Pattern (D ⊗ E)
\end{code}
%   ^^^cons   : (HIDE({D : U} →)) Pattern D → Pattern (list D) → Pattern (list D)
On all types we can use the |var| pattern which matches anything, while on sum types we can use the |left| and |right| patterns matching |inj₁| and |inj₂| constructors but not the |prod| pattern, which only makes sense for product types.
Also we have the |unit| pattern for matching |tt : ⊤|.

\begin{figure}
\beforefigure
\begin{code}
^^^pat-iso : (HIDE({D : U} ->)) (p : Pattern D) -> Iso ⟦ D ⟧ (PatResult p)
pat-iso p = record  {    to    = deconstruct p
                    SEP  from  = return ∘ construct p SEP ABBRV }
  where
    ^^^deconstruct :
      (HIDE({D : U})) (p : Pattern D) → ⟦ D ⟧ → Par (PatResult p)
    deconstruct var         x          = return x
    deconstruct unit        _          = return tt
    deconstruct (left   p)  (inj₁  x)  = deconstruct p x
    deconstruct (left   p)  (inj₂  y)  = fail
    deconstruct (right  p)  (inj₁  x)  = fail
    deconstruct (right  p)  (inj₂  y)  = deconstruct p y
    deconstruct (prod p q)  (x , y)    =  deconstruct p   x  >>= λ l  →
                                          deconstruct q   y  >>= λ r  →
                                          return (l , r)
    
    ^^^construct : (HIDE({D : U})) (p : Pattern D) → PatResult p → ⟦ D ⟧
    construct var         x         = x
    construct unit        _         = tt
    construct (left   p)  x         = inj₁ (construct p x)
    construct (right  p)  y         = inj₂ (construct p y)
    construct (prod p q)  (x , y)   = construct p x , construct q y
    
    ABBRV
\end{code}
\caption{Definition of the pattern matching isomorphism}
\label{fig:pat-iso}
\end{figure}

Next, patterns are given a different interpretation as ``containers'' for storing values --- whose types depend on the types of the variable sub-patterns --- at their variable positions.
For example, a pair pattern |prod var var : Pattern (D ⊗ E)| can be interpreted as a container storing a pair of type |f D × f E| for any given type-computing function |f : U -> Set|.
The types of such containers can be defined by induction on |Pattern|:
\begin{code}
^^^VarPositions :
  (HIDE({l : Level} {D : U} ->)) Pattern D -> (U -> Set (HIDE'(l))) -> Set (HIDE'(l))
VarPositions (var {D})   f {-"\kern-1pt"-}={-"\kern-1pt"-} f D
VarPositions unit        f {-"\kern-1pt"-}={-"\kern-1pt"-} ⊤
VarPositions (left   p)  f {-"\kern-1pt"-}={-"\kern-1pt"-} VarPositions p f
VarPositions (right  p)  f {-"\kern-1pt"-}={-"\kern-1pt"-} VarPositions p f
VarPositions (prod p q)  f {-"\kern-1pt"-}={-"\kern-1pt"-} VarPositions p f {-"\kern-1pt"-}×{-"\kern-1pt"-} VarPositions q f
\end{code}
A first situation where the notion of pattern-induced containers is helpful is when formalising pattern matching:
the result of a successful pattern matching can be filled into such a container, by putting the resulting components at their respective variable positions.
Defining the types of pattern matching results as an instance of |VarPositions|:
\begin{code}
^^^PatResult : (HIDE({D : U} ->)) Pattern D -> Set
PatResult p = VarPositions p ⟦_⟧
\end{code}
we can formulate pattern matching as a partial isomorphism:
\begin{code}
pat-iso : (HIDE({D : U} ->)) (p : Pattern D) -> Iso ⟦ D ⟧ (PatResult p)
\end{code}
whose definition is shown in \autoref{fig:pat-iso}.
The |to| component of the isomorphism performs pattern matching, and the |from| component reverses pattern matching by evaluating a pattern as if it is an expression, using the input of type |PatResult p| as an environment for values at variable positions.
(This pattern matching isomorphism will also play an important role in \autoref{sec:view-rearrangement}.)
%\footnote{
%For a concrete example: The pattern |pat = cons var var : Pattern (list A)| has two variable positions of types |⟦ A ⟧| and |List ⟦ A ⟧|, and indeed |PatResult pat| computes to |⟦ A ⟧ × List ⟦ A ⟧|.
%Thus |Iso.to (pat-iso pat)| matches any non-empty list and returns its head and tail as a pair.
%The pattern can also be viewed as an expression denoting a non-empty list; it is evaluated by applying |Iso.from (pat-iso pat)| to an environment of type |⟦ A ⟧ × List ⟦ A ⟧| supplying values for the head and the tail.}

For the source updating operation, we bypass formulation of a lens and directly define the semantics in terms of |interp| (due to some typing problem).
By specialising |VarPositions| we can place |BiGUL| statements --- along with their view types --- at the variable positions of a pattern:
\begin{code}
^^^BiGULs : (HIDE({D : U} ->)) Pattern D -> Set1
BiGULs p = VarPositions p (\ D -> (E : U) × BiGUL D E)
\end{code}
The |update| constructor of |BiGUL| then takes a pattern~|p| and a bunch of |BiGUL| statements collected in |bs : BiGULs p|.
Its view type is a product of the view types in~|bs|, whose code in~|U| is computed by
\begin{code}
^^^Views : (HIDE({D : U})) (p : Pattern D) -> BiGULs p -> U
Views var         (E , _)     = E
Views unit        _           = one
Views (left   p)  bs          = Views p bs
Views (right  p)  bs          = Views p bs
Views (prod p q)  (bs , bs')  = Views p bs ⊗ Views q bs'
\end{code}
The lens semantics of |update| is given by
\begin{code}
interp (update pat bs) =
  iso-lens (pat-iso pat) ↔ interp-update pat bs
\end{code}
where |interp-update| has type
\begin{code}
^^^interp-update :  (HIDE({D : U})) (p : Pattern D) (bs : BiGULs p) →
                    Lens (PatResult p) ⟦ Views p bs ⟧
\end{code}
which inductively interprets and composes the |BiGUL| statements in~|bs| into one lens.
In prose, the |put| semantics of |update| decomposes the source by matching it with the pattern~|p|, updates its components at the variable positions using the |BiGUL| statements in~|bs|, and reassembles the updated components, while its |get| semantics again decomposes the source by matching it with~|p| and extracts views from its components using~|bs|.

\subsection{View rearrangement}
\label{sec:view-rearrangement}

The operation |rearr| is for rearranging the view to a specific shape suitable for subsequent updates.
For example, the |update| operation in \autoref{sec:source-update} demands that the view must be in a shape that matches the source pattern (as computed by |Views|), and this is usually achieved by an outer |rearr|, as is the case in lines 3--4 of \autoref{fig:bookstore}.
In essence, |rearr| performs an invertible computation, while its |get| semantics performs the inverse of the computation.
The invertible computation used is usually just moving things around and not complicated (hence the name ``rearrangement''), so instead of letting the programmer write general invertible functions (probably along with their inverses), we ask instead for syntactic descriptions of simple invertible computations, which are more restrictive but allow automatic inversion.

The kind of invertible computation we wish to express is a decomposition of an old view followed by an assembling of a new view using all components of the old view --- for example, lines 3--4 of \autoref{fig:bookstore} simply express the computation:
\begin{code}
\ { (title , price) -> (title , tt , tt , price , tt) }
\end{code}
When it comes to decomposition, pattern matching should come to mind again; assembling can also be handled by pattern matching, as we have formulated pattern matching as an isomorphism |pat-iso| in \autoref{sec:source-update}.
We thus require two patterns |p : Pattern D| and |q : Pattern E|, the former for decomposing the old view of type |⟦ D ⟧| and the latter for assembling the new view of type |⟦ E ⟧|.
For the latter pattern~|q|, we also need to specify which component of the old view will be used as the value at each variable position.
This last piece of information is represented by specialising |VarPositions| again:
\begin{code}
^^^Paths : (HIDE({D E : U} ->)) Pattern D -> Pattern E -> Set1
Paths p q = VarPositions q (Path p)
\end{code}
where |Path p T|, defined below, is the type of all paths that lead from the root of a |PatResult p| to a component of type |⟦ T ⟧| at a variable position:
\begin{code}
^^^Path : (HIDE({D : U} →)) Pattern D → U → Set₁
Path (var {D})   T = D {-"\kern-2pt"-}≡{-"\kern-2pt"-} T
Path unit        T = ⊥
Path (left   p)  T = Path p T
Path (right  p)  T = Path p T
Path (prod p q)  T = Path p T ⊎ Path q T
\end{code}
%as evidenced by the following path evaluation function which extracts the component at the end of a path:
%\begin{code}
%^^^eval-path : {D T : U} (pat : Pattern D) → Path pat T → PatResult pat → ⟦ T ⟧
%\end{code}
If a given bunch of |paths : Paths p q| use up all components of the old view, i.e., all possible paths for~|p| are contained in |paths|, then |p|, |q|, and |paths| together specify a computation that can be automatically inverted --- since all components of the old view are present somewhere in the new view, we can always reassemble the old view from the new view (unless there is inconsistency --- a component of the old view might be copied into more than one positions of the new view, and the values at these positions of the new view must be the same when doing the reassembling).

In more detail, we can construct an isomorphism:
\begin{code}
^^^view-rearr-iso : (HIDE({D E : U} ->))
  (p : Pattern D) (q : Pattern E) (paths : Paths p q) →
  Invertible p q paths → Iso ⟦ E ⟧ ⟦ D ⟧
\end{code}
where |^^^Invertible p q paths| is defined to state that |paths| can pass a check for ensuring that all components of the old view are used.
The |from| direction matches an inhabitant of |⟦ D ⟧| with the pattern~|p| and uses the components as an environment for evaluating the pattern~|q| to an inhabitant of |⟦ E ⟧|.
The |to| direction is more delicate:
It creates an empty container of type |VarPositions p (Maybe ∘ ⟦_⟧)| (all of whose variable positions are |nothing|) and fills it with components obtained by matching the input |⟦ E ⟧| with~|q|.
If the container can be completely and consistently filled up --- that is, every position in the container is filled with one and only one value --- then we turn the container into a ``necessarily full'' container of type |VarPositions p ⟦_⟧| --- i.e., |PatResult p| --- by stripping all the |just| tags, and use it as an environment for evaluating~|p| to an inhabitant of |⟦ D ⟧|.

Finally, the lens semantics we assign to |rearr| is\todo{composition with iso}
\begin{code}
interp (rearr p q paths b) =
  interp b ↔ iso-lens (view-rearr-iso p q paths ABBRV)
\end{code}
This is not the actual definition, in fact --- note that the invertibility proof is not specified above.
To do that, we need to extend |interp| with an additional argument consisting of invertibility proofs for all |rearr| operations appearing in the program being interpreted.
This extension is omitted from the presentation for clarity.

\subsection{A showcase example: transatlantic corporation}
\label{sec:transatlantic-corporation}

We end this section by looking at an example in detail, which requires nontrivial putback logic, in particular involving |caseS| and |caseV|.
Suppose that a transatlantic corporation regularly relocates employees between its UK and US offices, and pay them in the local currency.
The payroll database stores for each employee their name, salary (a number interpreted as pounds or dollars depending on which country the employee works in), and current office location:
\begin{code}
^^^Employee = Name ⊗ (Salary ⊗ Location) : U
\end{code}
where |^^^Name : U| and |^^^Salary : U| are datatypes representing name and salary encoded as elements of~|U| (e.g., |k String| and |k ℕ|).
|Location| is defined as a disjoint sum:
\begin{code}
^^^Location = UKLocation ⊕ USLocation : U
\end{code}
where |^^^UKLocation : U| and |^^^USLocation : U| are, again, encoded datatypes representing UK and US office locations.
To manage relocation, we extract from the payroll database each employee's current office location:
\begin{code}
^^^Office = Name ⊗ Location : U
\end{code}
We would like to add, remove, or relocate employees by modifying the extracted list and put the modified list back to the payroll database.
Most interestingly, when an employee is relocated to a different country, their salary should also be converted to the local currency.
Below we describe this putback logic in |BiGUL|.

At top level, we align a list of |Employee|s with a list of |Office|s:
\begin{code}
^^^relocate-employees : BiGUL (list Employee) (list Office)
relocate-employees = align
  (\ _ -> return true)
  (\ { (sname , _) (vname , _) -> return (sname == vname) })
  relocate-employee
  (\ { (_ , loc) -> return ("" , 0 , loc) })
  (\ _ -> return nothing)
\end{code}
We are considering all |Employee|s in the list, so the source condition is always true.
Items from the two lists are matched by employee name.
Pairs of matched |Employee| and |Office| are synchronised by an inner |BiGUL| statement |relocate-employee|, to be defined later.
An unmatched |Office| means that a new employee is inserted, so we create a temporary |Employee| in the source list, which is then updated by |relocate-employee| using the unmatched |Office|.
(We put the view location in the temporary |Employee| merely for efficiency, as |relocate-employee| would not have to deal with mismatching source and view locations.)
An unmatched |Employee| means that the employee is removed from the view, so we conceal the item by deleting it from the source.

Matched pairs of |Employee| and |Office| are synchronised by
\begin{code}
^^^relocate-employee : BiGUL Employee Office
relocate-employee =
  rearr  (prod var var)  (prod  var (prod    unit  var))
                         (      inj₁ refl ,  tt ,  inj₂ refl)
         -- representing |\ { (name , loc) -> (name , tt , loc) }|
         (update  (prod   var              var)
                  ((_ ,   replace) , (_ ,  adjust-salary-and-office)))
\end{code}
First we rearrange the view to the same shape as the source --- specifically, we insert |tt| between the name and location, to pair with the salary in the source later.
Next we decompose the source by pattern matching to replace the name in the source with the name in the view --- this looks like a redundant step but is required because |BiGUL| enforces that all view information, in particular the name, must be put into the source.

The rest of the source, i.e., salary and location, is further synchronised with the remaining view by
\begin{code}
^^^adjust-salary-and-office : BiGUL  (Salary  ⊗ Location)
                                     (one     ⊗ Location)
adjust-salary-and-office = caseV
  (  (prod var (left   var) , relocate-to-UK  ) ∷
     (prod var (right  var) , relocate-to-US  ) ∷ [])
\end{code}
We match the view location with either the |left| or |right| pattern to know which country the employee will work in.
If the employee will work in the UK, the update to the source is
\begin{code}
^^^relocate-to-UK : BiGUL  (Salary  ⊗ Location    )
                           (one     ⊗ UKLocation  )
relocate-to-UK = caseS
  (  (inUK  , normal'    (update  (prod  var (left     var))
                                  ((_ ,  skip) , (_ ,  replace)))) ∷
     (inUS  , adaptive'  (λ {  (salary , _) →
                               {-"~"-} return (toPounds salary , inj₁ "") })) ∷ [])
\end{code}
We probe which country the employee is in originally with a |caseS|.
The predicates |^^^inUK| and |^^^inUS| have type
\begin{code}
⟦ Salary ⟧ × (⟦ UKLocation ⟧ ⊎ ⟦ USLocation ⟧) → Par Bool
\end{code}
and are simply defined as
\begin{code}
inUK  (_ , inj₁  _) = return true
inUK  (_ , inj₂  _) = return false

inUS  (_ , inj₁  _) = return false
inUS  (_ , inj₂  _) = return true
\end{code}
If the original country is the UK, the employee is not being relocated to a different country, so we can skip updating the salary and just replace the location.
If it is the US, we adapt the source by changing the salary from dollars to pounds (by a function |^^^toPounds : Salary -> Salary| defined elsewhere) and change the location to an empty one in the UK, which will be replaced by the view location in the second round of the source case analysis following the adaptation.
The other branch |relocate-to-US| is defined symmetrically:
\begin{code}
^^^relocate-to-US : BiGUL  (Salary  ⊗ Location    )
                           (one     ⊗ USLocation  )
relocate-to-US = caseS
  (  (inUS  , normal'    (update  (prod  var (right    var))
                                  ((_ ,  skip) , (_ ,  replace)))) ∷
     (inUK  , adaptive'  (λ {  (salary , _) →
                               {-"~"-} return (£⇒$ salary , inj₂ "") })) ∷ [])
\end{code}

We should now check\todo{Hugo: why do we need to check? Only for the purpose of ensuring totality, I suppose.}\ whether all constraints for passing dynamic checking are met, working from the inner statements towards the outer ones.
For the |caseS| statements in |relocate-to-UK| and |relocate-to-US|, the normal branches do not change the source country and thus do not switch branch, and the adaptive branches do change the source so it matches a normal branch.
For the |caseV| in |adjust-salary-and-office|, the ranges of |relocate-to-UK| and |relocate-to-US| are disjoint since they respectively produce UK and US locations.
For the |align| in |relocate-employees|, since the source condition is always true, whatever |relocate-employee| produces necessarily satisfies the source condition, and we must conceal unmatched sources by deleting them, which is indeed what we specify; finally, we must guarantee that the updated source produced by |relocate-employee| and the view again satisfy the matching condition, even when the view is an unmatched one and the original source is a temporary one created from the view, and this is indeed guaranteed because |relocate-employee| puts the name in the view into the source.


\section{Concluding remarks}
\label{sec:concluding-remarks}

We have given an account of the design and semantics of |BiGUL| and also how the language and its well-behavedness are formalised and proved in \Agda.
Below we conclude this paper by commenting on the putback-based approach to bidirectional transformations, emphasising our formal development, discussing issues about dynamic checking, comparing |BiGUL| with its closest relative \Putlenses, and talking briefly about porting |BiGUL| to \Haskell.

One might wonder how |BiGUL|, being described as a putback-based language, stands out from existing work on lenses ---
after all, the fundamental definition of |Lens| is exactly the classic symmetric definition.\todo{Hugo: classical lenses are asymmetric}\
We designed |BiGUL| such that the natural way to understand a |BiGUL| program is to read it as a description of putback logic, and, more importantly, a |BiGUL| program can be written by thinking solely in the putback direction (as opposed to programming the |get| direction or, worse, switching between |get| and |put|).
(It may appear that the programmer needs to devise a |get| before programming a corresponding |put|, but that impression most likely stems from the confusion between general BX specifications --- e.g., consistency \emph{relations} between sources and views --- and |get| \emph{functions} as a much more specialised form of BX specifications.
The programmer should of course have a specification in mind before programming |put|, but that specification is not necessarily a |get| function.)

We admit, though, that the expressive power of the current |BiGUL| is not obviously stronger than existing lenses (although the adaptive |put| semantics of |caseS| has shown some potential of putback-based design), and it might be hard to find some formal properties to differentiate the putback-based approach in this paper from other approaches to lenses.
We thus believe that the main contribution of this paper should be its completely formal development:
Grohne~et~al~\cite{Grohne-formalizing-semantic-bidirectionalization} have also worked on formal verification about BXs (in \Agda), but their work was about semantic bidirectionalisation based on parametric polymorphism~\cite{Voigtlander-BX-for-free}.
Our work is the first to formally describe and verify a collection of nontrivial and practically implemented lenses in a dependently typed language, in particular utilising the power of dependent types to design a logically clean abstract syntax and make well-behavedness proofs much easier to handle.

%\todo[inline]{It is probably worth noting that the well-behavedness proofs in |align-lens| are more difficult, but it is mostly due to the inherent difficulty of the proofs --- there is little overhead regarding the formalisation per~se, due to the nice structure of the definition of well-behavedness (in terms of~|_↦_|).
%The exact details of the aligning algorithm are pretty delicate (especially the last part for merging the sources that satisfy |source-condition| with those that do not), and these details being formally verified makes us absolutely confident that the algorithm is correct.}

One might be sceptical about the use of dynamic checks to guarantee the well-behavedness of |BiGUL|, as the programmer can write programs that fail inadvertently and only realise that at runtime.
We currently deal with this problem by informally stating the constraints that should be satisfied by the programmer-supplied actions in order to pass the dynamic checks.
It is also possible to turn such statements into formal theorems saying that a |put| succeeds if relevant constraints are satisfied.
We do not take this step because, as future work, we plan to jump further by switching to a dependently typed setting where constraints on programmer-supplied actions can be directly specified in their types. This will not only eliminate all dynamic checks from the lens semantics, but also help the programmer to write correct putback programs with the type system (as demonstrated by, e.g., Norell~\cite{Norell-interactive-programming}).
%(This is one of the reasons that we choose \Agda\ as the formalisation language: since \Agda\ is particularly suited for dependently typed programming (as opposed to verification of simply typed programs using dependent types), using \Agda\ would save some effort doing the transition.)

|BiGUL| is closely related to \Putlenses~\cite{Pacheco-putlenses}:
Both |BiGUL| and \Putlenses\ are datatype-generic and describe partial putback transformations that use dynamic checks to guarantee well-behavedness.
|BiGUL| does not strive for maximal expressiveness like \Putlenses\ (in particular, |BiGUL| does not consider effectful computations in |put|), but instead takes a more cautious approach by focusing on a smaller set of essential features and formally proving their well-behavedness.
The dynamic checks are sometimes tricky to get right, and proofs about them can require detailed case analyses that are not so easy to manage.
(Indeed, during the formalisation we did not get everything right the first time and had to make corrections after failing to complete the proofs.)
Thus the fact that the whole |BiGUL| language is formally verified is a firm step towards reliable putback-based bidirectional programming: as |BiGUL| is intended to serve as the (new) core of higher-level putback-based languages like \BiFluX~\cite{Pacheco-BiFluX} and \BiYacc~\cite{Zhu-BiYacc}, the well-behavedness of the latter languages will be established beyond doubt.

To (re-)implement \BiFluX\ and \BiYacc, we have ported |BiGUL| to \Haskell.
The types need some adaptation because \Haskell\ is not fully dependently typed (but close enough; see, e.g., Lindley and McBride~\cite{Lindley-Hasochism}), while the transformations can be more or less faithfully transcribed.
A significant difference between \Agda\ and \Haskell\ is that \Haskell\ freely allows general recursion, whereas \Agda\ requires every recursive program to be evidently well-founded.
We do not consider recursion in the \Agda\ formalisation, but we need general recursion in \BiFluX\ and \BiYacc, which is a major reason for porting |BiGUL| to \Haskell.
We believe that the total correctness proved in \Agda\ can be translated into some sort of partial correctness in \Haskell\ (perhaps along the line of Danielsson~et~al~\cite{Danielsson-fast-and-loose-reasoning}); that is, a |BiGUL| program in \Haskell\ may not terminate, but when it does, it is guaranteed to be well-behaved.

\section*{Acknowledgements}

We would like to thank Li Liu, Hugo Pacheco, and the anonymous reviewers of APLAS and PEPM for their valuable comments.

%\section{Conclusion}
%\label{sec:con}
%% reification ?
%While putback-based bidirectional programming opens a new research area,
%it also brings up a challenge of ensuring well-behavedness of putback transformations.
%In this paper, we design and implement a completely formally verified language |BiGUL|
%that can support development of robust putback transformations.
%We do not design a programmer-friendly syntax for |BiGUL|, as it is supposed to serve as the core language 
%for other general putback-based BX languages such as \BiFluX~\cite{Pacheco-BiFluX} and \BiYacc~\cite{Zhu-BiYacc}.
%Different from the core language for
%putback-based bidirectional programming, it has two 
%distinguishing operations, the align operation and the extended source case analysis 
%operation with adaption capability. We have successfully given a proof-carrying semantics
%of |BiGUL| using the dependently typed programming language \Agda,
%showing that a |get| can be derived automatically from |put|, and that the pair of |get| and |put| 
%does form a well-behaved bidirectional transformation. 

% programming.
%which aims to make BX to be deployed for really practical usages.
%Thus not only allow specifying flexibility over the updates like existing putback-based languages,
%but also need fully guarantee the correctness of derived $get$.
%\BiGUL\ is a datatype-generic, and fully certified BX core language
%which can be served as the base for other putback-based languages,
%and it is indeed that BiFLuX can be translated into \BiGUL.
%\todo[inline]{Towards a dependently typed version}

%Currently alignment only handles the case that
%elements with the same matching key in both source and view are unique
%%, which indicates the update operations after alignment focus on single source/view element.
%, but in fact there are situations that more than one elements have the same matching key value,
%the result would be a list of elements for both source and view.

%\acks
%
%Acknowledgments, if needed.

%\appendix
%\section{Appendix Title}
%
%This is the text of the appendix, if you need one.

\bibliographystyle{abbrvnat}

\bibliography{ref}

\end{document}
